
\documentclass[10pt,openany,a4paper]{article}
\usepackage{tikz, pgfplots, tkz-euclide,calc}
\usetikzlibrary{patterns,shapes.arrows,calc}
\usepackage{fancyhdr}
\usepackage{enumerate,enumitem}
\usepackage{cancel}
\usepackage{varwidth}

% TAMBAHKAN PACKAGE SENDIRI KALAU KURANG

\usepackage{geometry}
\geometry{
	left = 22mm,
	right = 22mm,
	top = 25mm,
	bottom = 25mm,
}
\usepackage{graphicx}
\usepackage{subcaption}
\usepackage{soul}
\usepackage{hyperref}
\usepackage{float}
\usepackage{setspace}
\usepackage{afterpage}
\usepackage{xcolor}
\usepackage[nosectionbib]{apacite}
\usepackage{lmodern}
\usepackage{mathtools}
\usepackage{upgreek}
\usepackage{siunitx}
% \usepackage{fontawesome5}
\usepackage{fontawesome6}
\usepackage{physics}
\usepackage{pdfpages}
\usepackage{hyperref}
\renewcommand{\baselinestretch}{1.3}
\usepackage{multicol}
\usepackage{hyperref,ragged2e,amsmath,multicol,gensymb,setspace,
fancyhdr,amsfonts,tikz,pgfplots,nccmath,enumerate,verbatim,amsthm,physics}
\usepgfplotslibrary{polar,fillbetween}
% Disable tikz externalization to avoid requiring -shell-escape
\usepgfplotslibrary{external}
\tikzexternalize[prefix=tikz/]
\tikzset{reuse path/.code={\pgfsyssoftpath@setcurrentpath{#1}}}
\usetikzlibrary{shapes.geometric,arrows}
\newcommand\mylog[1]{\mathop{{}^{#1}\mathrm{log}}}
\pgfplotsset{compat=1.15}
\pgfplotsset{my style/.append style={axis x line=middle, axis y line=
middle, xlabel={$x$}, ylabel={$y$}, axis equal }}
\usepackage{etoolbox}
% remove stray \makeatother
\newtheorem{defn}{Definisi}[section]
\newtheorem{teo}[defn]{Teorema}
\newtheorem{thm}{Teorema}[section]
\newtheorem{lemma}{Lemma}
\newtheorem{lemmas}[defn]{Lemma}
\newtheorem{cor}[defn]{Akibat}
\newtheorem{asumsi}[defn]{Asumsi}
\theoremstyle{definition}
\newtheorem{con}{Contoh}[section]
\theoremstyle{definition}
\theoremstyle{plain}
\newtheorem{prop}{Proposisi}[section]
\renewcommand{\proofname}{Bukti}
\renewcommand{\thethm}{\arabic{chapter}.\arabic{thm}}

\newcommand{\normlp}[1]{\left\|#1\right\|_{L^p_{\lambda}(0,T;H)}}% Fungsi norm (||x||)
\newcommand{\R}{\mathbb{R}}
\renewcommand{\natural}{\mathbb{N}}
\newcommand{\Xpl}{X^R_{p,\lambda}}
\newcommand{\Xpls}{X^R_{p,\lambda^*}}
\newcommand{\qpl}{q^R_{p,\lambda}}
\newcommand{\qpls}{q^R_{p,\lambda^*}}
\pagestyle{fancy}
\fancyhf{}
\lhead{Halaman \thepage}
\rhead{Kumpulan Soal ETS Kalkulus 1}
% \cfoot{\faXTwitter\ \href{https://x.com/ahmadzakiyudin_/}{@ahmadzakiyudin\_} | \faInstagram\ \href{https://instagram.com/ahmadzakiyudin_/}{@ahmadzakiyudin\_} | \faTiktok\ \href{https://tiktok.com/@ahmathz/}{@ahmathz}}
% \rhead{Pembahasan Soal ETS \\ Kalkulus 1 ITS 2025/2026}
\hypersetup{
    colorlinks=true,
    linkcolor=blue,
    filecolor=blue,      
    urlcolor=blue,
}
\newenvironment{jawab}[2][tes]{
% \fbox{\begin{minipage}{14.65cm}
\textbf{Kata kunci materi:} #1.\\
\textbf{Penyelesaian:} #2
% \end{minipage}}
}{}

\newenvironment{jawab2}[2][tes]{
\fbox{\begin{minipage}{14.65cm}
 #2
\end{minipage}}}{}

\newenvironment{jawaban}[2][tes]{
\colorlet{oldcolor}{.}
        \color{red}
        \setlength{\fboxrule}{1pt}\boxed{\color{oldcolor}#2}\color{oldcolor}.
}{}

\newenvironment{jawaban2}[2][tes]{
\colorlet{oldcolor}{.}
        \color{red}
        \setlength{\fboxrule}{1pt}\fbox{\color{oldcolor}#2}\color{oldcolor}.
}{}

\newenvironment{penting}[2][tes]{
\colorlet{oldcolor}{.}
\color{blue}
\fbox{\color{oldcolor}#2}\color{oldcolor}
}{}

\usepackage{multirow}
\begin{document}
% \begin{titlepage}
%     \vspace*{\fill}
%     \begin{center}
%       \Huge {PEMBAHASAN SOAL EAS \\ KALKULUS II \\ TAHUN 2024/2025}\\[0.4 cm]
%       \huge {Ahmad Hisbu Zakiyudin}\\[0.2 cm]
%       \Large {(\href{https://instagram.com/ahmadzakiyudin_/}{@ahmadzakiyudin\_})}
%     \end{center}
%     \vspace*{\fill}
%   \end{titlepage}
% !TEX root = ../../BookTemplate.tex
%%%%%%%%%%%%%%%%%%%%%%%%%%%%%%%%%%%%%%%%%%%%%%%%%%%%%%%%%%%%%%%%%%%%%%%%%%%%%%%%%%
% \begin{titlepage}
%     \raggedleft	
% 	\hspace{.025\textwidth}
% 	\parbox[b]{\textwidth}{
%     \vspace{1.5cm}
% 		{\Huge\bfseries Pembahasan Soal EAS}                \\[20pt]
% 		{\huge\textit\ Kalkulus II}\\ \\[10pt]
%         {\Large\textsc{Tahun 2024/2025}}}
% \vfill
% \rule{1pt}{.14\textheight}
% \hspace{.025\textwidth}
%     \raisebox{10pt}{\parbox[b]{.95\textwidth}{ 
%         {\Large\textsc{Ahmad Hisbu Zakiyudin}}              \\[10pt]
%         \large{\href{https://instagram.com/ahmadzakiyudin_/}{@ahmadzakiyudin\_}}               \\[25pt]
%         {\large 5 Juni 2026}}}
% \end{titlepage}
\begin{titlepage}
	\centering % Center everything on the title page
	% \scshape % Use small caps for all text on the title page
	\vspace*{1.5\baselineskip} % White space at the top of the page
% ===================
%	Title Section 	
% ===================

	\rule{13cm}{1.6pt}\vspace*{-\baselineskip}\vspace*{2pt} % Thick horizontal rule
	\rule{13cm}{0.4pt} % Thin horizontal rule
	
		\vspace{0.75\baselineskip} % Whitespace above the title
% ========== Title ===============	
	{	\Huge Kumpulan Soal ETS\\ 
			\vspace{4mm}
		Kalkulus 1 \\ \vspace{2mm} \huge{Tahun 2023/2024 -- 2025/2026}\\ \vspace{2mm} \Large{(Disusun per Bab dan Subbab)}\\	}
% ======================================
		\vspace{0.75\baselineskip} % Whitespace below the title
	\rule{13cm}{0.4pt}\vspace*{-\baselineskip}\vspace{3.2pt} % Thin horizontal rule
	\rule{13cm}{1.6pt} % Thick horizontal rule
	
		\vspace{1.75\baselineskip} % Whitespace after the title block
% =================
%	Information	
% =================
{\large Diketik ulang oleh: Ahmad Hisbu Zakiyudin}
	% {\large Oleh: Ahmad Hisbu Zakiyudin \\
	% 	\vspace*{1.2\baselineskip}
	% \href{https://instagram.com/ahmadzakiyudin_/}{@ahmadzakiyudin\_}} \\
	\vfill
	\vspace{1mm}
\huge{Selamat Belajar!}
\end{titlepage}
\makeatletter
\setlength\fboxsep{3pt}
% ensure enough header height for fancyhdr
\setlength{\headheight}{24pt}

\setlength{\headheight}{24pt}

\section*{Keterangan}
Kumpulan ini memuat seluruh soal ETS Kalkulus 1 ITS (SM234101) Tahun 2023/2024, 2024/2025, dan 2025/2026, dikelompokkan berdasarkan bab dan subbab pada buku ajar Kalkulus 1. Label \textbf{[Subbab \dots]} di sebelah nomor menunjukkan semua subbab yang dicakup oleh soal tersebut; soal ditempatkan pada subbab utamanya. Keterangan sumber (tahun, hari, kelas, dan nomor soal asli) dituliskan di bawah label. Soal kelas IUP hari Selasa 2025/2026 identik dengan soal kelas 113 dan 114, sehingga hanya ditulis sekali.

\begin{center}
\begin{tabular}{|c|l|c|}
\hline
\textbf{Bab} & \textbf{Subbab} & \textbf{Banyak soal}\\ \hline
\multirow{4}{*}{1} & 1.1 Bilangan Real & 9\\ \cline{2-3}
 & 1.2 Nilai Mutlak & 7\\ \cline{2-3}
 & 1.4 Persamaan Garis dan Jarak & 7\\ \cline{2-3}
 & 1.5 Persamaan Lingkaran dan Parabola & 3\\ \hline
\multirow{3}{*}{2} & 2.2 Operasi pada Fungsi & 18\\ \cline{2-3}
 & 2.3 Grafik Fungsi & 4\\ \cline{2-3}
 & 2.4 Fungsi Invers & 15\\ \hline
\multirow{4}{*}{3} & 3.1 Pengantar Notasi Limit & 1\\ \cline{2-3}
 & 3.2 Penghitungan Limit & 6\\ \cline{2-3}
 & 3.3 Limit di Tak-Hingga & 3\\ \cline{2-3}
 & 3.4 Kekontinuan & 10\\ \hline
\multirow{3}{*}{4} & 4.2 Fungsi Turunan & 1\\ \cline{2-3}
 & 4.3 Diferensiasi & 3\\ \cline{2-3}
 & 4.4 Aturan Rantai dan Diferensiasi Implisit & 18\\ \hline
\multicolumn{2}{|r|}{\textbf{Total}} & \textbf{105}\\ \hline
\end{tabular}
\end{center}
\newpage

\section*{Bab 1 \quad Sistem Bilangan Real}
\addcontentsline{toc}{section}{Bab 1 Sistem Bilangan Real}
\subsection*{1.1 \quad Bilangan Real}
\begin{enumerate}
\setcounter{enumi}{0}
    \item \textbf{[Subbab 1.1]} {\footnotesize\textcolor{gray}{(ETS 2023/2024 -- Senin, 16 Oktober 2023, Kelas 11--16, No.~1)}}\par
    Dapatkan himpunan penyelesaian dari pertidaksamaan
    \begin{align*}
        \frac{\frac{1}{3}x-5}{2x-5}\geq 3
    \end{align*}

    \item \textbf{[Subbab 1.1]} {\footnotesize\textcolor{gray}{(ETS 2023/2024 -- Senin, 16 Oktober 2023, Kelas 31--37, No.~1)}}\par
    Dapatkan himpunan penyelesaian dari pertidaksamaan
    \begin{align*}
        \frac{x-1}{2-x}\geq\frac{x+1}{x+2}
    \end{align*}

    \item \textbf{[Subbab 1.1]} {\footnotesize\textcolor{gray}{(ETS 2023/2024 -- Selasa, 17 Oktober 2023, Kelas 17--23, No.~1)}}\par
    Dapatkan himpunan penyelesaian pertidaksamaan
    \begin{align*}
        2\leq\frac{x^2+1}{x}<x+3
    \end{align*}

    \item \textbf{[Subbab 1.1]} {\footnotesize\textcolor{gray}{(ETS 2024/2025 -- Kamis, 17 Oktober 2024, Kelas 5--12, 101, No.~1)}}\par
    Dapatkan himpunan penyelesaian dari $\dfrac{1}{x+2}<\dfrac{1}{4-x}$.

    \item \textbf{[Subbab 1.1]} {\footnotesize\textcolor{gray}{(ETS 2025/2026 -- Senin, 13 Oktober 2025, Kelas 40--49, 62--66, No.~1)}}\par
    Dapatkan semua nilai $x$ yang memenuhi pertidaksamaan
    \[
        \dfrac{2x + 1}{3x - 6} < \dfrac{1}{2}.
    \]

    \item \textbf{[Subbab 1.1]} {\footnotesize\textcolor{gray}{(ETS 2025/2026 -- Monday, 13 October 2025, Class 1, 4, 67, 70, 71 (IUP), No.~1)}}\par
    Solve the following inequality.
    \[ \frac{x^2 - x - 12}{x^2 - 2x - 3} \ge 0. \]

    \item \textbf{[Subbab 1.1]} {\footnotesize\textcolor{gray}{(ETS 2025/2026 -- Selasa, 14 Oktober 2025, Kelas 113, 114 (juga kelas IUP 2, 3, 68, 69), No.~1)}}\par
    Dapatkan penyelesaian dari pertidaksamaan
    \begin{align*}
        \frac{x^2-x-12}{x^2+3x-10}\geq 0.
    \end{align*}

    \item \textbf{[Subbab 1.1]} {\footnotesize\textcolor{gray}{(ETS 2025/2026 -- Selasa, 14 Oktober 2025, Kelas 5--12, 23--29, 115, 116, No.~1)}}\par
    Selesaikan pertidaksamaan berikut
    \begin{align*}
        \frac{1}{x+2} \leq \frac{2}{x-1}.
    \end{align*}

    \item \textbf{[Subbab 1.1]} {\footnotesize\textcolor{gray}{(ETS 2025/2026 -- Selasa, 14 Oktober 2025, Kelas 30--39, No.~1)}}\par
    Dapatkan penyelesaian dari pertidaksamaan
    \begin{align*}
        2\leq x^2-2\leq 7.
    \end{align*}

\end{enumerate}
\subsection*{1.2 \quad Nilai Mutlak}
\begin{enumerate}
\setcounter{enumi}{9}
    \item \textbf{[Subbab 1.2]} {\footnotesize\textcolor{gray}{(ETS 2023/2024 -- Selasa, 17 Oktober 2023, Kelas 24--30, No.~1)}}\par
    Tentukan semua nilai $x$ yang memenuhi
    \begin{align*}
        1<\frac{1}{|x-2|}\leq 4.
    \end{align*}

    \item \textbf{[Subbab 1.2]} {\footnotesize\textcolor{gray}{(ETS 2023/2024 -- Selasa, 17 Oktober 2023, Kelas 38--43, 60, No.~1)}}\par
    Dapatkan penyelesaian dari $|2x+1|=|3x-2|+1$.

    \item \textbf{[Subbab 1.2]} {\footnotesize\textcolor{gray}{(ETS 2024/2025 -- Kamis, 17 Oktober 2024, Kelas 47--59, 111, No.~1)}}\par
    Dapatkan himpunan penyelesaian dari $\dfrac{1}{|2x-5|}>5$.

    \item \textbf{[Subbab 1.2]} {\footnotesize\textcolor{gray}{(ETS 2024/2025 -- Kamis, 17 Oktober 2024, Kelas 20--33, 105, 106, No.~1)}}\par
    Dapatkan himpunan penyelesaian dari $-1\leq |2-x|<3$.

    \item \textbf{[Subbab 1.2]} {\footnotesize\textcolor{gray}{(ETS 2025/2026 -- Senin, 13 Oktober 2025, Kelas 17--22, No.~1)}}\par
    \begin{enumerate}
        \item Nyatakan pernyataan $\sqrt{(2x-1)^2}$ dalam bentuk nilai mutlak.
        \item Dapatkan semua nilai $x$ yang memenuhi persamaan $\sqrt{(2x-1)^2}=1-2x$.
    \end{enumerate}

    \item \textbf{[Subbab 1.2]} {\footnotesize\textcolor{gray}{(ETS 2025/2026 -- Senin, 13 Oktober 2025, Kelas 13--16, 50--61, No.~1)}}\par
    \begin{enumerate}
        \item Definisikan $|3-x|$.
        \item Berdasarkan point (a) tersebut, tentukan semua nilai $x$ yang memenuhi pertidaksamaan
        \begin{align*}
            |3-x|\leq -2x+1.
        \end{align*}
    \end{enumerate}

    \item \textbf{[Subbab 1.2]} {\footnotesize\textcolor{gray}{(ETS 2025/2026 -- Selasa, 14 Oktober 2025, Kelas 102, 108--112, No.~1)}}\par
    Dapatkan solusi dari ketaksamaan
    \begin{align*}
        2<|3-x|\leq 5.
    \end{align*}
    Petunjuk: Gunakan rumus jarak.

\end{enumerate}
\subsection*{1.4 \quad Persamaan Garis dan Jarak}
\begin{enumerate}
\setcounter{enumi}{16}
    \item \textbf{[Subbab 1.4]} {\footnotesize\textcolor{gray}{(ETS 2023/2024 -- Senin, 16 Oktober 2023, Kelas 4--10, No.~1)}}\par
    Diberikan tiga titik koordinat $A(-1,3)$, $B(1,-2)$ dan $C(3,0)$.
    \begin{enumerate}
        \item Dapatkan persamaan garis lurus $\ell$ yang melalui titik $B$ dan $C$.
        \item Dapatkan persamaan garis lurus $k$ yang melalui titik $A$ dan tegak lurus dengan garis $\ell$.
        \item Dapatkan jarak titik $A$ ke garis lurus $\ell$.
    \end{enumerate}

    \item \textbf{[Subbab 1.4]} {\footnotesize\textcolor{gray}{(ETS 2023/2024 -- Selasa, 17 Oktober 2023, Kelas 44--50, No.~1)}}\par
    Persamaan garis $\ell$ melalui titik $A(2,1)$ dan $B(-2,3)$.
    \begin{enumerate}
        \item Dapatkan persamaan garis $\ell$.
        \item Dapatkan persamaan garis $k$ yang tegak lurus dengan garis $\ell$ yang melalui titik tengah antara $A$ dan $B$.
        \item Sketsa grafik persamaan garis $\ell$ dan $k$ dalam satu bidang koordinat.
    \end{enumerate}

    \item \textbf{[Subbab 1.4]} {\footnotesize\textcolor{gray}{(ETS 2023/2024 -- Selasa, 17 Oktober 2023, Kelas 51--56, 61, No.~1)}}\par
    Diberikan garis $\ell$ melalui titik $(6,4)$ dan $(a,b)$ dengan kemiringan $2$, dan garis $k$ melalui titik $(0,0)$ dan $(a,b)$ dengan kemiringan $\frac{1}{2}$. Dapatkan nilai $a$ dan $b$.

    \item \textbf{[Subbab 1.4]} {\footnotesize\textcolor{gray}{(ETS 2024/2025 -- Rabu, 16 Oktober 2024, Kelas 13--19, 103, No.~1)}}\par
    Diberikan titik $A(2,-1)$, $B(2,2)$ dan $C(0,4)$. Dapatkan persamaan garis yang melalui titik $A$ dan sejajar dengan garis yang melalui $B$ dan $C$.

    \item \textbf{[Subbab 1.4]} {\footnotesize\textcolor{gray}{(ETS 2024/2025 -- Kamis, 17 Oktober 2024, Kelas 34--46, 107, 108, No.~1)}}\par
    Diberikan titik $A(1,1)$, $B(4,2)$ dan $C(2,6)$. Tentukan jarak dari titik $C$ ke garis yang melalui titik $A$ dan $B$.

    \item \textbf{[Subbab 1.4]} {\footnotesize\textcolor{gray}{(ETS 2025/2026 -- Senin, 13 Oktober 2025, Kelas 40--49, 62--66, No.~5)}}\par
    Diketahui garis $g_1: y = -x + 6$ dan garis $g_2: y = 2x$.
    \begin{enumerate}
        \item Tentukan titik potong kedua garis $g_1$ dan $g_2$.
        \item Tentukan titik potong $g_1$ dengan garis $x = 1$.
        \item Tentukan pula titik potong $g_2$ dengan garis $x = 1$. Gambarkan sketsa ketiga garis $g_1$, $g_2$, dan garis $x = 1$. Hitung luas segitiga yang dibentuk ketiga garis tersebut.
    \end{enumerate}

    \item \textbf{[Subbab 1.4]} {\footnotesize\textcolor{gray}{(ETS 2025/2026 -- Selasa, 14 Oktober 2025, Kelas 102, 108--112, No.~2)}}\par
    Diketahui garis $\ell_1$ dengan gradien $m_1=\frac{2}{3}$ dan garis $\ell_2$ berpotongan tegak lurus dengan $\ell_1$ di titik $(4,-1)$. Dapatkan persamaan garis $\ell_2$. \\
    Gambarkan kedua garis $\ell_1$ dan $\ell_2$ pada satu bidang koordinat.

\end{enumerate}
\subsection*{1.5 \quad Persamaan Lingkaran dan Parabola}
\begin{enumerate}
\setcounter{enumi}{23}
    \item \textbf{[Subbab 1.5]} {\footnotesize\textcolor{gray}{(ETS 2025/2026 -- Senin, 13 Oktober 2025, Kelas 17--22, No.~3)}}\par
    Diberikan persamaan separuh lingkaran
    \begin{align*}
       x=4+\sqrt{-y^2+14y-24}.
    \end{align*}
    \begin{enumerate}
        \item Dapatkan pusat dan jari-jari lingkaran.
        \item Sketsa grafiknya.
    \end{enumerate}

    \item \textbf{[Subbab 1.4, 1.5]} {\footnotesize\textcolor{gray}{(ETS 2025/2026 -- Senin, 13 Oktober 2025, Kelas 13--16, 50--61, No.~2)}}\par
    \begin{enumerate}
        \item Tentukan jarak titik $(1,3)$ dan $(7,11)$.
        \item Tentukan persamaan lingkaran yang diameternya adalah jarak titik $(1,3)$ ke $(7,11)$.
    \end{enumerate}

    \item \textbf{[Subbab 1.4, 1.5]} {\footnotesize\textcolor{gray}{(ETS 2025/2026 -- Selasa, 14 Oktober 2025, Kelas 113, 114 (juga kelas IUP 2, 3, 68, 69), No.~2)}}\par
    \begin{enumerate}
        \item Dapatkan kemiringan garis $\ell$ yang melalui titik $(0,0)$ dan $(-2,1)$.
        \item Jika garis $\ell$ tegak lurus dengan garis yang menyinggung lingkaran yang berpusat di $(0,0)$ dan berjari-jari $\sqrt{5}$, di titik $(-2,1)$, maka dapatkan persamaan garis singgung lingkaran tersebut dan sketsalah.
    \end{enumerate}

\end{enumerate}
\newpage

\section*{Bab 2 \quad Fungsi}
\addcontentsline{toc}{section}{Bab 2 Fungsi}
\subsection*{2.2 \quad Operasi pada Fungsi}
\begin{enumerate}
\setcounter{enumi}{0}
    \item \textbf{[Subbab 2.1, 2.2]} {\footnotesize\textcolor{gray}{(ETS 2023/2024 -- Senin, 16 Oktober 2023, Kelas 4--10, No.~2)}}\par
    Diketahui $f(x)=\dfrac{4}{x^2-1}$ dan $g(x)=\sqrt{x+4}$.
    \begin{enumerate}
        \item Dapatkan domain $f$ dan $g$.
        \item Dapatkan $(g\circ f)(x)$ beserta domainnya.
    \end{enumerate}

    \item \textbf{[Subbab 2.1, 2.2]} {\footnotesize\textcolor{gray}{(ETS 2023/2024 -- Senin, 16 Oktober 2023, Kelas 11--16, No.~2)}}\par
    Diberikan $f(x)=x^2-1$ dan $g(x)=\sqrt{x-3}$.
    \begin{enumerate}
        \item Dapatkan domain $f$ dan $g$.
        \item Dapatkan $(g\circ f)(x)$ beserta domainnya.
    \end{enumerate}

    \item \textbf{[Subbab 2.1, 2.2]} {\footnotesize\textcolor{gray}{(ETS 2023/2024 -- Senin, 16 Oktober 2023, Kelas 31--37, No.~2)}}\par
    Diketahui $f(x)=\sqrt{x+1}$ dan $g(x)=\sqrt{4-x^2}$.
    \begin{enumerate}
        \item Dapatkan domain $f$ dan $g$.
        \item Dapatkan $(g\circ f)(x)$ beserta domainnya.
    \end{enumerate}

    \item \textbf{[Subbab 2.1, 2.2]} {\footnotesize\textcolor{gray}{(ETS 2023/2024 -- Selasa, 17 Oktober 2023, Kelas 17--23, No.~2)}}\par
    Diberikan fungsi $f(x)=\sqrt{5-x}$ dan $g(x)=\dfrac{1}{9-x^2}$.
    \begin{enumerate}
        \item Dapatkan domain $f$ dan $g$.
        \item Dapatkan $(g\circ f)(x)$ beserta domainnya.
    \end{enumerate}

    \item \textbf{[Subbab 2.1, 2.2]} {\footnotesize\textcolor{gray}{(ETS 2023/2024 -- Selasa, 17 Oktober 2023, Kelas 24--30, No.~2)}}\par
    Diketahui $f(x)=\dfrac{1+x}{1-x}$ dan $g(x)=\dfrac{x}{1-x}$.
    \begin{enumerate}
        \item Dapatkan domain $f$ dan $g$.
        \item Dapatkan $(f\circ g)(x)$ beserta domainnya.
    \end{enumerate}

    \item \textbf{[Subbab 2.1, 2.2]} {\footnotesize\textcolor{gray}{(ETS 2023/2024 -- Selasa, 17 Oktober 2023, Kelas 38--43, 60, No.~2)}}\par
    Diberikan fungsi $f(x)=|x^2-1|$ dan $g(x)=\dfrac{1}{x^2-9}$.
    \begin{enumerate}
        \item Dapatkan domain $f$ dan $g$.
        \item Dapatkan $(g\circ f)(x)$ beserta domainnya.
    \end{enumerate}

    \item \textbf{[Subbab 2.1, 2.2]} {\footnotesize\textcolor{gray}{(ETS 2023/2024 -- Selasa, 17 Oktober 2023, Kelas 44--50, No.~2)}}\par
    Diberikan fungsi $f(x)=\sqrt{1-x}$ dan $g(x)=\dfrac{1}{\sqrt{9-x^2}}$.
    \begin{enumerate}
        \item Dapatkan domain $f$ dan $g$.
        \item Dapatkan $(g\circ f)(x)$ beserta domainnya.
    \end{enumerate}

    \item \textbf{[Subbab 2.1, 2.2]} {\footnotesize\textcolor{gray}{(ETS 2023/2024 -- Selasa, 17 Oktober 2023, Kelas 51--56, 61, No.~2)}}\par
    Diberikan fungsi $f(x)=-\sqrt{|x^2-16|}$ dan $g(x)=\dfrac{1}{x^2-9}$.
    \begin{enumerate}
        \item Dapatkan domain $f$ dan $g$.
        \item Dapatkan $(g\circ f)(x)$ beserta domainnya.
    \end{enumerate}

    \item \textbf{[Subbab 2.1, 2.2]} {\footnotesize\textcolor{gray}{(ETS 2024/2025 -- Rabu, 16 Oktober 2024, Kelas 13--19, 103, No.~2)}}\par
    Diberikan $f(x)=\dfrac{1}{x^2-4}$ dan $g(x)=\sqrt{x+1}$.
    \begin{enumerate}
        \item Dapatkan domain $f(x)$ dan $g(x)$.
        \item Dapatkan $(f\circ g)(x)$ dan domain $(f\circ g)(x)$.
    \end{enumerate}

    \item \textbf{[Subbab 2.1, 2.2]} {\footnotesize\textcolor{gray}{(ETS 2024/2025 -- Kamis, 17 Oktober 2024, Kelas 5--12, 101, No.~2)}}\par
    Diberikan $f(x)=x^2+2,\ x\geq 0$ dan $g(x)=\sqrt{x-3}$.
    \begin{enumerate}
        \item Dapatkan domain $f(x)$ dan $g(x)$.
        \item Dapatkan $(g\circ f)(x)$ dan domain $(g\circ f)(x)$.
    \end{enumerate}

    \item \textbf{[Subbab 2.1, 2.2]} {\footnotesize\textcolor{gray}{(ETS 2024/2025 -- Kamis, 17 Oktober 2024, Kelas 47--59, 111, No.~2)}}\par
    Diberikan $f(x)=\sqrt{x-3}$ dan $g(x)=1+\sqrt{x-5}$.
    \begin{enumerate}
        \item Dapatkan domain $f(x)$ dan $g(x)$.
        \item Dapatkan $(f\circ g)(x)$ dan domain $(f\circ g)(x)$.
    \end{enumerate}

    \item \textbf{[Subbab 2.1, 2.2]} {\footnotesize\textcolor{gray}{(ETS 2024/2025 -- Kamis, 17 Oktober 2024, Kelas 20--33, 105, 106, No.~2)}}\par
    Diberikan $f(x)=\sqrt{25-x^2}$ dan $g(x)=\dfrac{1}{x^2}$.
    \begin{enumerate}
        \item Dapatkan domain $f(x)$ dan $g(x)$.
        \item Dapatkan $(g\circ f)(x)$ dan domain $(g\circ f)(x)$.
    \end{enumerate}

    \item \textbf{[Subbab 2.1, 2.2]} {\footnotesize\textcolor{gray}{(ETS 2024/2025 -- Kamis, 17 Oktober 2024, Kelas 34--46, 107, 108, No.~2)}}\par
    Diberikan $f(x)=x-4$, untuk $x\geq 4$ dan $g(x)=\sqrt{4-x^2}$.
    \begin{enumerate}
        \item Dapatkan domain $f(x)$ dan $g(x)$.
        \item Tentukan $(g\circ f)(x)$ dan domain $(g\circ f)(x)$.
    \end{enumerate}

    \item \textbf{[Subbab 2.1, 2.2]} {\footnotesize\textcolor{gray}{(ETS 2025/2026 -- Senin, 13 Oktober 2025, Kelas 40--49, 62--66, No.~2)}}\par
    Diberikan fungsi $f(x) = \dfrac{1}{x-1}$ dan $g(x) = 2 - \sqrt{x}$.
    \begin{enumerate}
        \item Tentukan domain dari fungsi $f$ dan $g$.
        \item Dapatkan fungsi $(f \circ g)(x)$ beserta domainnya.
    \end{enumerate}

    \item \textbf{[Subbab 2.1, 2.2]} {\footnotesize\textcolor{gray}{(ETS 2025/2026 -- Monday, 13 October 2025, Class 1, 4, 67, 70, 71 (IUP), No.~2)}}\par
    Given a function $f(x) = 2 - \sqrt{x}$ and $g(x) = \dfrac{1}{x-1}$. Determine:
    \begin{enumerate}
        \item[a)] The domain of function $f$ and $g$;
        \item[b)] Function $(f \circ g)(x)$ and its domain.
    \end{enumerate}

    \item \textbf{[Subbab 2.2]} {\footnotesize\textcolor{gray}{(ETS 2025/2026 -- Selasa, 14 Oktober 2025, Kelas 113, 114 (juga kelas IUP 2, 3, 68, 69), No.~5)}}\par
    Diberikan fungsi $f(x)=\dfrac{1}{x-1}$ dan $g(x)=x^2$.
    \begin{enumerate}
        \item Dapatkan $(f\circ g)(x)$.
        \item Jika diketahui $(p\circ g)(x)=x^2-3$, dapatkan $p(x)$.
    \end{enumerate}

    \item \textbf{[Subbab 2.1, 2.2]} {\footnotesize\textcolor{gray}{(ETS 2025/2026 -- Selasa, 14 Oktober 2025, Kelas 5--12, 23--29, 115, 116, No.~3)}}\par
    Diberikan fungsi $f(x)=\dfrac{1}{x-2}$ dan $g(x)=\sqrt{5-x}$. Tentukan:
    \begin{enumerate}
        \item Domain fungsi $f$ dan $g$.
        \item Fungsi $(f\circ g)(x)$ dan $\left(\dfrac{f}{g}\right) (x)$ beserta domainnya masing-masing.
    \end{enumerate}

    \item \textbf{[Subbab 2.2]} {\footnotesize\textcolor{gray}{(ETS 2025/2026 -- Selasa, 14 Oktober 2025, Kelas 30--39, No.~3)}}\par
    Tentukan domain dari fungsi $(f\circ g)(x)$ jika diketahui fungsi $g(x)=2x+3$ dengan $0\leq x<3$ dan fungsi $f$ terdefinisi pada interval $-1\leq x <5$.

\end{enumerate}
\subsection*{2.3 \quad Grafik Fungsi}
\begin{enumerate}
\setcounter{enumi}{18}
    \item \textbf{[Subbab 1.5, 2.3]} {\footnotesize\textcolor{gray}{(ETS 2025/2026 -- Monday, 13 October 2025, Class 1, 4, 67, 70, 71 (IUP), No.~5)}}\par
    Sketch the base curve $x = y^2$. Based on that base curve, use translation to sketch the curve $x = y^2 + 4y + 1$. (Explain the process).

    \item \textbf{[Subbab 2.1, 2.3]} {\footnotesize\textcolor{gray}{(ETS 2025/2026 -- Senin, 13 Oktober 2025, Kelas 13--16, 50--61, No.~3)}}\par
    Diberikan fungsi $f(x)=\begin{cases}
        x-1, &x<3\\
        2(x-4)^2, &x\geq 3
    \end{cases}$.
    \begin{enumerate}
        \item Dapatkan $f(3)$, $f(-2)$, dan $f(1-x^2)$.
        \item Dapatkan domain dan range dari $f(x)$.
        \item Sketsa grafik dari $f(x)$.
    \end{enumerate}

    \item \textbf{[Subbab 2.3]} {\footnotesize\textcolor{gray}{(ETS 2025/2026 -- Selasa, 14 Oktober 2025, Kelas 5--12, 23--29, 115, 116, No.~2)}}\par
    Gambarkan kurva dasar $y=\dfrac{1}{x}$. Berdasarkan kurva dasar tersebut, gunakan translasi (pergeseran) untuk menggambar kurva $y=\dfrac{1}{x-2}+4$ (Jelaskan ilustrasinya).

    \item \textbf{[Subbab 1.5, 2.3]} {\footnotesize\textcolor{gray}{(ETS 2025/2026 -- Selasa, 14 Oktober 2025, Kelas 30--39, No.~2)}}\par
    Gunakan grafik dasar $x=y^2$ dan sifat pergeseran/translasi untuk menggambar parabola $x=y^2+2y+3$.

\end{enumerate}
\subsection*{2.4 \quad Fungsi Invers}
\begin{enumerate}
\setcounter{enumi}{22}
    \item \textbf{[Subbab 2.4]} {\footnotesize\textcolor{gray}{(ETS 2023/2024 -- Senin, 16 Oktober 2023, Kelas 4--10, No.~3)}}\par
    Diberikan fungsi $f(x)=1+\sqrt{8-x^3}$.
    \begin{enumerate}
        \item Dapatkan invers fungsi $f(x)$.
        \item Dapatkan domain dan range invers fungsinya.
    \end{enumerate}

    \item \textbf{[Subbab 2.4]} {\footnotesize\textcolor{gray}{(ETS 2023/2024 -- Senin, 16 Oktober 2023, Kelas 11--16, No.~3)}}\par
    Diberikan fungsi $f(x)=x^2-4x+5$.
    \begin{enumerate}
        \item Tentukan nilai $k$ terkecil sehingga $f(x)$ memiliki invers pada $[k,\infty)$.
        \item Tentukan invers fungsi $f(x)$.
    \end{enumerate}

    \item \textbf{[Subbab 2.3, 2.4]} {\footnotesize\textcolor{gray}{(ETS 2023/2024 -- Senin, 16 Oktober 2023, Kelas 31--37, No.~3)}}\par
    Diberikan suatu fungsi $f(x)=2+\sqrt{4-x^2}$ dengan domain $f$ adalah $[a,2]$.
    \begin{enumerate}
        \item Tentukan nilai terkecil $a$ agar $f^{-1}(x)$ ada.
        \item Dapatkan $f^{-1}(x)$.
        \item Sketsa grafik $f(x)$ dan $f^{-1}(x)$ pada satu bidang koordinat.
    \end{enumerate}

    \item \textbf{[Subbab 2.1, 2.4]} {\footnotesize\textcolor{gray}{(ETS 2023/2024 -- Selasa, 17 Oktober 2023, Kelas 17--23, No.~3)}}\par
    Diketahui $f(x)=\sqrt{4-x^2}$.
    \begin{enumerate}
        \item Tentukan domain $f(x)$.
        \item Apakah $f^{-1}(x)$ ada? Jelaskan alasannya.
        \item Jika tidak ada, definisikan domain $f(x)$ agar $f^{-1}(x)$ ada.
    \end{enumerate}

    \item \textbf{[Subbab 2.3, 2.4]} {\footnotesize\textcolor{gray}{(ETS 2023/2024 -- Selasa, 17 Oktober 2023, Kelas 24--30, No.~3)}}\par
    Diberikan fungsi $f(x)=x^2-2x+1$.
    \begin{enumerate}
        \item Tentukan nilai $k$ terbesar sehingga $f(x)$ memiliki invers pada $(-\infty,k]$.
        \item Tentukan $f^{-1}(x)$.
        \item Sketsa grafik $f(x)$ dan $f^{-1}(x)$ dalam satu bidang koordinat.
    \end{enumerate}

    \item \textbf{[Subbab 2.1, 2.4]} {\footnotesize\textcolor{gray}{(ETS 2023/2024 -- Selasa, 17 Oktober 2023, Kelas 38--43, 60, No.~3)}}\par
    Diberikan fungsi sebagai berikut:
    \begin{align*}
        f(x)=-\frac{2}{x+3};\quad \text{dan}\quad g(x)=\frac{3x+2}{x+2}
    \end{align*}
    \begin{enumerate}
        \item Tentukan domain dan range dari masing-masing fungsi tersebut.
        \item Apakah kedua fungsi tersebut saling invers? Jelaskan alasannya.
    \end{enumerate}

    \item \textbf{[Subbab 2.4]} {\footnotesize\textcolor{gray}{(ETS 2023/2024 -- Selasa, 17 Oktober 2023, Kelas 44--50, No.~3)}}\par
    Diberikan fungsi $f(x)=\sqrt[3]{x+1}+2$.
    \begin{enumerate}
        \item Dapatkan $f^{-1}(x)$.
        \item Dapatkan domain dan range dari $f^{-1}(x)$.
    \end{enumerate}

    \item \textbf{[Subbab 2.3, 2.4]} {\footnotesize\textcolor{gray}{(ETS 2023/2024 -- Selasa, 17 Oktober 2023, Kelas 51--56, 61, No.~3)}}\par
    Diberikan suatu fungsi $f(x)=2-\sqrt{4x-x^2}$ dengan domain $f$ adalah $[0,a]$.
    \begin{enumerate}
        \item Dapatkan nilai terbesar $a$ sehingga $f^{-1}(x)$ ada.
        \item Tentukan $f^{-1}(x)$.
        \item Gambarkan sketsa kurva $f(x)$ dan $f^{-1}(x)$ dalam satu bidang koordinat.
    \end{enumerate}

    \item \textbf{[Subbab 2.4]} {\footnotesize\textcolor{gray}{(ETS 2024/2025 -- Rabu, 16 Oktober 2024, Kelas 13--19, 103, No.~3)}}\par
    Diberikan $f(x)=x^2-4x+7$.
    \begin{enumerate}
        \item Tentukan domain dari $f$ sehingga $f^{-1}$ ada.
        \item Dapatkan $f^{-1}$ beserta domainnya.
    \end{enumerate}

    \item \textbf{[Subbab 2.3, 2.4]} {\footnotesize\textcolor{gray}{(ETS 2024/2025 -- Kamis, 17 Oktober 2024, Kelas 5--12, 101, No.~3)}}\par
    Diketahui $f(x)=x^3-2$.
    \begin{enumerate}
        \item Dapatkan $f^{-1}(x)$ beserta domainnya.
        \item Sketsa grafik dari $f(x)$ dan $f^{-1}(x)$ pada satu bidang koordinat.
    \end{enumerate}

    \item \textbf{[Subbab 2.3, 2.4]} {\footnotesize\textcolor{gray}{(ETS 2024/2025 -- Kamis, 17 Oktober 2024, Kelas 47--59, 111, No.~3)}}\par
    Diberikan $f(x)=\sqrt[3]{x-1},\ x\geq 1$.
    \begin{enumerate}
        \item Dapatkan $f^{-1}(x)$ beserta domainnya.
        \item Sketsa $f(x)$ dan $f^{-1}(x)$ pada satu bidang koordinat.
    \end{enumerate}

    \item \textbf{[Subbab 2.4]} {\footnotesize\textcolor{gray}{(ETS 2024/2025 -- Kamis, 17 Oktober 2024, Kelas 20--33, 105, 106, No.~3)}}\par
    Diberikan $f(x)=\sqrt{4-x^2}$.
    \begin{enumerate}
        \item Tentukan domain dari $f(x)$ sehingga inversnya ada.
        \item Dapatkan $f^{-1}(x)$ beserta domainnya.
    \end{enumerate}

    \item \textbf{[Subbab 2.3, 2.4]} {\footnotesize\textcolor{gray}{(ETS 2024/2025 -- Kamis, 17 Oktober 2024, Kelas 34--46, 107, 108, No.~3)}}\par
    Diberikan fungsi $f(x)=\sqrt{2x-x^2},\ 0\leq x\leq 1$.
    \begin{enumerate}
        \item Tentukan $f^{-1}(x)$ beserta domainnya.
        \item Gambarkan grafik $f(x)$ dan $f^{-1}(x)$ dalam satu bidang koordinat.
    \end{enumerate}

    \item \textbf{[Subbab 2.1, 2.4]} {\footnotesize\textcolor{gray}{(ETS 2025/2026 -- Senin, 13 Oktober 2025, Kelas 17--22, No.~2)}}\par
    Diberikan fungsi $f(x)=\sqrt{x-4}$.
    \begin{enumerate}
        \item Tentukan domain dan range dari fungsi $f(x)$.
        \item Tentukan $f^{-1}(x)$, domain dan range dari $f^{-1}(x)$.
    \end{enumerate}

    \item \textbf{[Subbab 2.4]} {\footnotesize\textcolor{gray}{(ETS 2025/2026 -- Selasa, 14 Oktober 2025, Kelas 102, 108--112, No.~3)}}\par
    Diketahui $f(x)=x^2-4x$.
    \begin{enumerate}
        \item Tentukan domain alami untuk $x$ agar $f(x)$ mempunyai invers.
        \item Dapatkan $f^{-1}(x)$, domain, dan rangenya.
    \end{enumerate}

\end{enumerate}
\newpage

\section*{Bab 3 \quad Limit dan Kekontinuan}
\addcontentsline{toc}{section}{Bab 3 Limit dan Kekontinuan}
\subsection*{3.1 \quad Pengantar Notasi Limit}
\begin{enumerate}
\setcounter{enumi}{0}
    \item \textbf{[Subbab 3.1]} {\footnotesize\textcolor{gray}{(ETS 2023/2024 -- Senin, 16 Oktober 2023, Kelas 4--10, No.~4)}}\par
    Diberikan $f(x)=\begin{cases}
        2k^2(x-1), & x\leq 2\\
        \left(\dfrac{1-k}{2}\right)x, & x>2
    \end{cases}$.\\
    Tentukan semua nilai $k$ yang memenuhi agar $\displaystyle\lim_{x\to 2}f(x)$ ada.

\end{enumerate}
\subsection*{3.2 \quad Penghitungan Limit}
\begin{enumerate}
\setcounter{enumi}{1}
    \item \textbf{[Subbab 3.2]} {\footnotesize\textcolor{gray}{(ETS 2023/2024 -- Selasa, 17 Oktober 2023, Kelas 17--23, No.~4)}}\par
    Hitunglah $\displaystyle\lim_{x\to 3}\frac{\sqrt{x^2-5}-2}{x-3}$.

    \item \textbf{[Subbab 3.2]} {\footnotesize\textcolor{gray}{(ETS 2024/2025 -- Kamis, 17 Oktober 2024, Kelas 34--46, 107, 108, No.~4)}}\par
    Hitunglah $\displaystyle\lim_{x\to 2}\frac{\sqrt{x^2+5}-3}{x-2}$.

    \item \textbf{[Subbab 3.2, 3.4]} {\footnotesize\textcolor{gray}{(ETS 2025/2026 -- Senin, 13 Oktober 2025, Kelas 17--22, No.~5)}}\par
    Diberikan $f(x)=\begin{cases}
        \dfrac{\sqrt{x+9}-3}{x}, &x\neq 0\\[2mm]
        \dfrac{1}{3}k-\dfrac{1}{5}, &x=0
    \end{cases}$, dengan $x\in [-9,\infty)$.\\
    Tentukan nilai $k$ sehingga $\displaystyle\lim_{x\to 0}f(x)=f(0)$.\\
    Kemudian nyatakan fungsi $f(x)$ tanpa menggunakan bentuk fungsi sepotong-sepotong dalam domainnya untuk nilai $k$ yang diperoleh.

    \item \textbf{[Subbab 2.3, 3.2]} {\footnotesize\textcolor{gray}{(ETS 2025/2026 -- Senin, 13 Oktober 2025, Kelas 13--16, 50--61, No.~4)}}\par
    Diberikan $f(x)=\dfrac{x^3-1}{x-1}$.
    \begin{enumerate}
        \item Tentukan $\displaystyle\lim_{x\to 1}f(x)$.
        \item Buatlah sketsa grafik $y=f(x)$.
    \end{enumerate}

    \item \textbf{[Subbab 3.2]} {\footnotesize\textcolor{gray}{(ETS 2025/2026 -- Selasa, 14 Oktober 2025, Kelas 113, 114 (juga kelas IUP 2, 3, 68, 69), No.~4)}}\par
    Dapatkan $\displaystyle \lim_{x\to 3} \frac{\sqrt{3x+7}-4}{x-3}$.

    \item \textbf{[Subbab 3.2, 3.4]} {\footnotesize\textcolor{gray}{(ETS 2025/2026 -- Selasa, 14 Oktober 2025, Kelas 102, 108--112, No.~4)}}\par
    Diberikan fungsi
    \begin{align*}
        f(x)=\begin{cases}
            \dfrac{x^2-9}{x+3}, &x\neq -3\\
            k, &x=-3
        \end{cases}
    \end{align*}
    \begin{enumerate}
        \item Tentukan nilai $k$ agar $f(-3)=\displaystyle\lim_{x\to -3} f(x)$.
        \item Dengan memberi $k$ pada $\displaystyle\lim_{x\to-3} f(x)$, tunjukkan bahwa $f(x)$ dapat dinyatakan sebagai fungsi polinomial.
    \end{enumerate}

\end{enumerate}
\subsection*{3.3 \quad Limit di Tak-Hingga}
\begin{enumerate}
\setcounter{enumi}{7}
    \item \textbf{[Subbab 3.3]} {\footnotesize\textcolor{gray}{(ETS 2023/2024 -- Selasa, 17 Oktober 2023, Kelas 38--43, 60, No.~4)}}\par
    Hitunglah $\displaystyle\lim_{x\to-\infty}\frac{5-3x}{\sqrt{7+27x^2}}$.

    \item \textbf{[Subbab 3.3]} {\footnotesize\textcolor{gray}{(ETS 2024/2025 -- Kamis, 17 Oktober 2024, Kelas 5--12, 101, No.~4)}}\par
    Hitunglah $\displaystyle\lim_{x\to-\infty}\frac{\sqrt{5x^2-2}}{x+3}$.

    \item \textbf{[Subbab 3.3]} {\footnotesize\textcolor{gray}{(ETS 2024/2025 -- Kamis, 17 Oktober 2024, Kelas 47--59, 111, No.~4)}}\par
    Hitunglah $\displaystyle\lim_{y\to-\infty}\frac{2-y}{\sqrt{7+4y^2}}$.

\end{enumerate}
\subsection*{3.4 \quad Kekontinuan}
\begin{enumerate}
\setcounter{enumi}{10}
    \item \textbf{[Subbab 3.1, 3.4]} {\footnotesize\textcolor{gray}{(ETS 2023/2024 -- Senin, 16 Oktober 2023, Kelas 11--16, No.~4)}}\par
    Diketahui $f(x)=\dfrac{1}{|3-x|}$.
    \begin{enumerate}
        \item Hitunglah $\displaystyle\lim_{x\to 3}f(x)$.
        \item Berapa nilai $f(3)$?
        \item Apakah $f(x)$ kontinu? jelaskan!
        \item Apakah diskontinuitas $f(x)$ dapat dihilangkan? jika iya, definisikan kembali $f(x)$ sehingga menjadi kontinu.
    \end{enumerate}

    \item \textbf{[Subbab 3.2, 3.4]} {\footnotesize\textcolor{gray}{(ETS 2023/2024 -- Senin, 16 Oktober 2023, Kelas 31--37, No.~4)}}\par
    Diberikan $f(x)=\begin{cases}
        \dfrac{x^3-x^2+x-1}{x-1}, & x\neq 1\\
        k, & x=1
    \end{cases}$.\\
    Tentukan nilai $k$, sehingga $f(x)$ kontinu.

    \item \textbf{[Subbab 3.1, 3.4]} {\footnotesize\textcolor{gray}{(ETS 2023/2024 -- Selasa, 17 Oktober 2023, Kelas 24--30, No.~4)}}\par
    Diberikan fungsi $f(x)=\begin{cases}
        2|x+1|, & x<-1\\
        x^2, & x\geq -1
    \end{cases}$.
    \begin{enumerate}
        \item Hitunglah $\displaystyle\lim_{x\to-1}f(x)$.
        \item Berapa nilai $f(-1)$?
        \item Apakah $f(x)$ kontinu di $x=-1$? Jelaskan!
    \end{enumerate}

    \item \textbf{[Subbab 3.1, 3.4]} {\footnotesize\textcolor{gray}{(ETS 2023/2024 -- Selasa, 17 Oktober 2023, Kelas 44--50, No.~4)}}\par
    Diketahui $f(x)=\begin{cases}
        3x+1, & x>2\\
        -2x+11, & x<2
    \end{cases}$
    \begin{enumerate}
        \item Hitunglah $\displaystyle\lim_{x\to 2}f(x)$.
        \item Berapa $f(2)$?
        \item Apakah terdapat diskontinuitas dari $f(x)$? Jelaskan alasannya.
        \item Jika terdapat diskontinuitas, apakah dapat dihilangkan? Jika iya, definisikan kembali $f(x)$ agar kontinu.
    \end{enumerate}

    \item \textbf{[Subbab 1.2, 3.4]} {\footnotesize\textcolor{gray}{(ETS 2023/2024 -- Selasa, 17 Oktober 2023, Kelas 51--56, 61, No.~4)}}\par
    Diberikan fungsi $f(x)=\dfrac{16-x^2}{4-|x|}$.
    \begin{enumerate}
        \item Nyatakan fungsi $f(x)$ dalam bentuk fungsi sepotong-sepotong.
        \item Selidiki apakah $f(x)$ memiliki titik diskontinu.
        \item Jika memiliki titik diskontinu, apakah titik diskontinu dapat dihilangkan?
    \end{enumerate}

    \item \textbf{[Subbab 3.1, 3.4]} {\footnotesize\textcolor{gray}{(ETS 2024/2025 -- Rabu, 16 Oktober 2024, Kelas 13--19, 103, No.~4)}}\par
    Dapatkan nilai $k$ sedemikian sehingga fungsi:
    \begin{align*}
        f(x)=\begin{cases}
            x^2-k, & x<3\\
            3x-3, & x\geq 3
        \end{cases}
    \end{align*}
    kontinu di $x=3$.

    \item \textbf{[Subbab 3.2, 3.4]} {\footnotesize\textcolor{gray}{(ETS 2024/2025 -- Kamis, 17 Oktober 2024, Kelas 20--33, 105, 106, No.~4)}}\par
    Diberikan fungsi
    \begin{align*}
        f(x)=\begin{cases}
            \dfrac{4-x}{2-\sqrt{x}}, & x\neq 4\\
            6, & x=4
        \end{cases},
    \end{align*}
    Selidiki kekontinuan $f(x)$ di $x=4$.

    \item \textbf{[Subbab 3.2, 3.4]} {\footnotesize\textcolor{gray}{(ETS 2025/2026 -- Senin, 13 Oktober 2025, Kelas 40--49, 62--66, No.~3)}}\par
    Diberikan fungsi $f(x) = \begin{cases} \dfrac{x^3-8}{2x-4}, & x \neq 2 \\ k, & x=2 \end{cases}$.
    \begin{enumerate}
        \item Tentukan nilai $k$ sehingga $f(x)$ kontinu di $x = 2$.
        \item Berdasarkan hasil dari (a), nyatakan $f(x)$ sebagai polinomial.
    \end{enumerate}

    \item \textbf{[Subbab 3.2, 3.4]} {\footnotesize\textcolor{gray}{(ETS 2025/2026 -- Monday, 13 October 2025, Class 1, 4, 67, 70, 71 (IUP), No.~4)}}\par
    Given $f(x) = \begin{cases} \dfrac{x^3-8}{2x-4}, & x \neq 2 \\ 2k, & x=2 \end{cases}$.
    \begin{enumerate}
        \item[a.] Determine $k$ such that $f(x)$ continuous at $x = 2$.
        \item[b.] Based on the result of a), rewrite $f(x)$ as a polynomial.
    \end{enumerate}

    \item \textbf{[Subbab 1.2, 2.3, 3.4]} {\footnotesize\textcolor{gray}{(ETS 2025/2026 -- Selasa, 14 Oktober 2025, Kelas 30--39, No.~4)}}\par
    Diberikan fungsi $f(x)=\left|\dfrac{x^2-9}{x-3}\right|$. Dapatkan titik $x$ di mana membuat fungsi $f(x)$ diskontinu?\\
    Dapatkan nilai fungsi $f(x)$ pada titik tersebut agar menjadi kontinu dan sketsa grafik fungsi tersebut.

\end{enumerate}
\newpage

\section*{Bab 4 \quad Derivatif}
\addcontentsline{toc}{section}{Bab 4 Derivatif}
\subsection*{4.2 \quad Fungsi Turunan}
\begin{enumerate}
\setcounter{enumi}{0}
    \item \textbf{[Subbab 3.4, 4.2]} {\footnotesize\textcolor{gray}{(ETS 2025/2026 -- Selasa, 14 Oktober 2025, Kelas 5--12, 23--29, 115, 116, No.~4)}}\par
    Tunjukkan bahwa $f(x)=\begin{cases}
        3x+1, &x\leq 2\\
        x+5, &x>2
    \end{cases}$ adalah fungsi kontinu tetapi tidak dapat diturunkan di $x=2$.

\end{enumerate}
\subsection*{4.3 \quad Diferensiasi}
\begin{enumerate}
\setcounter{enumi}{1}
    \item \textbf{[Subbab 1.4, 4.3]} {\footnotesize\textcolor{gray}{(ETS 2023/2024 -- Selasa, 17 Oktober 2023, Kelas 17--23, No.~5)}}\par
    Diberikan $f(x)=\dfrac{x+3}{x+2}$.
    \begin{enumerate}
        \item Dapatkan $\dfrac{dy}{dx}$.
        \item Jika $(x_0,y_0)$ merupakan titik pada kurva $f(x)$, maka dapatkan titik $(x_0,y_0)$ dan garis singgung di $(x_0,y_0)$ yang tegak lurus dengan $y=x$.
    \end{enumerate}

    \item \textbf{[Subbab 4.3]} {\footnotesize\textcolor{gray}{(ETS 2025/2026 -- Monday, 13 October 2025, Class 1, 4, 67, 70, 71 (IUP), No.~3)}}\par
    Write the derivative formula of $\dfrac{f}{g}$, that is $\dfrac{d}{dx}\left(\dfrac{f(x)}{g(x)}\right)$. \\
    Then, find the derivative of $\dfrac{f}{g}$ given that $f(x) = 2x^2 + 3x - 2$ and $g(x) = x - 1$ for $x \neq 1$.

    \item \textbf{[Subbab 4.3]} {\footnotesize\textcolor{gray}{(ETS 2025/2026 -- Selasa, 14 Oktober 2025, Kelas 113, 114 (juga kelas IUP 2, 3, 68, 69), No.~3)}}\par
    Tuliskan rumus turunan dari $fg$ yaitu $\dfrac{d}{dx}(f(x)g(x))$.\\
    Selanjutnya, dapatkan turunan dari $fg$, jika diketahui $f(x)=3x^2-7x$ dan $g(x)=3-2x^2$.

\end{enumerate}
\subsection*{4.4 \quad Aturan Rantai dan Diferensiasi Implisit}
\begin{enumerate}
\setcounter{enumi}{4}
    \item \textbf{[Subbab 4.4]} {\footnotesize\textcolor{gray}{(ETS 2023/2024 -- Senin, 16 Oktober 2023, Kelas 4--10, No.~5)}}\par
    Gunakan diferensiasi implisit untuk mendapatkan $\dfrac{dy}{dx}$ dari $y-x^2y^3=y^2x-x$.

    \item \textbf{[Subbab 4.4]} {\footnotesize\textcolor{gray}{(ETS 2023/2024 -- Senin, 16 Oktober 2023, Kelas 11--16, No.~5)}}\par
    Gunakan diferensiasi implisit untuk mendapatkan $\dfrac{dy}{dx}$ dari $x^2y^3-2y^2-3x+7=0$.

    \item \textbf{[Subbab 4.4]} {\footnotesize\textcolor{gray}{(ETS 2023/2024 -- Senin, 16 Oktober 2023, Kelas 31--37, No.~5)}}\par
    Dapatkan persamaan garis singgung kurva $x^2y^2+4xy=12y$ di titik $(2,1)$.

    \item \textbf{[Subbab 1.4, 4.4]} {\footnotesize\textcolor{gray}{(ETS 2023/2024 -- Selasa, 17 Oktober 2023, Kelas 24--30, No.~5)}}\par
    Diketahui garis singgung pada kurva $y^3=2x^2$ di titik $(a,b)$ sejajar dengan garis $2x+3y-5=0$. Tentukan titik $(a,b)$.

    \item \textbf{[Subbab 4.4]} {\footnotesize\textcolor{gray}{(ETS 2023/2024 -- Selasa, 17 Oktober 2023, Kelas 38--43, 60, No.~5)}}\par
    Dapatkan persamaan garis singgung dari kurva $x^3-2x^2y+2xy^2-y^3+3=0$ yang melalui titik $(1,2)$.

    \item \textbf{[Subbab 4.3, 4.4]} {\footnotesize\textcolor{gray}{(ETS 2023/2024 -- Selasa, 17 Oktober 2023, Kelas 44--50, No.~5)}}\par
    Gunakan diferensiasi implisit untuk mendapatkan $\dfrac{dy}{dx}$ dari $x^2-y+\dfrac{y}{x}=5$.

    \item \textbf{[Subbab 4.3, 4.4]} {\footnotesize\textcolor{gray}{(ETS 2023/2024 -- Selasa, 17 Oktober 2023, Kelas 51--56, 61, No.~5)}}\par
    Dapatkan $\dfrac{dy}{dx}$ dari fungsi $y=2x^2+\sqrt{\dfrac{x-3}{x^2+1}}$.

    \item \textbf{[Subbab 4.3, 4.4]} {\footnotesize\textcolor{gray}{(ETS 2024/2025 -- Rabu, 16 Oktober 2024, Kelas 13--19, 103, No.~5)}}\par
    Dapatkan $f'(x)$ dimana $f(x)=\sqrt{\dfrac{(3x+1)^3}{2x}}$.

    \item \textbf{[Subbab 4.4]} {\footnotesize\textcolor{gray}{(ETS 2024/2025 -- Kamis, 17 Oktober 2024, Kelas 5--12, 101, No.~5)}}\par
    Dapatkan persamaan garis singgung kurva $xy^2+y+\sqrt{x}=x+3$ di titik $(4,1)$.

    \item \textbf{[Subbab 4.3, 4.4]} {\footnotesize\textcolor{gray}{(ETS 2024/2025 -- Kamis, 17 Oktober 2024, Kelas 47--59, 111, No.~5)}}\par
    Dapatkan persamaan garis singgung kurva $x^2+y+\dfrac{y}{x}=\sqrt{x}+2$ di titik $(1,1)$.

    \item \textbf{[Subbab 4.3, 4.4]} {\footnotesize\textcolor{gray}{(ETS 2024/2025 -- Kamis, 17 Oktober 2024, Kelas 20--33, 105, 106, No.~5)}}\par
    Dapatkan $f''(x)$ dimana $f(x)=2x+\left(2\sqrt{x}-3\right)^{-2}$.

    \item \textbf{[Subbab 1.4, 4.4]} {\footnotesize\textcolor{gray}{(ETS 2024/2025 -- Kamis, 17 Oktober 2024, Kelas 34--46, 107, 108, No.~5)}}\par
    Diketahui persamaan garis singgung kurva $y^2-xy=-2$ di titik $(a,b)$ sejajar dengan kurva $y=-x$. Dapatkan titik $(a,b)$.

    \item \textbf{[Subbab 4.4]} {\footnotesize\textcolor{gray}{(ETS 2025/2026 -- Senin, 13 Oktober 2025, Kelas 40--49, 62--66, No.~4)}}\par
    Diberikan fungsi implisit $3x^2 + y^2 - 2xy + y - x = 6$. Tentukan $\dfrac{dy}{dx}$ di titik $(0,-3)$.

    \item \textbf{[Subbab 4.4]} {\footnotesize\textcolor{gray}{(ETS 2025/2026 -- Senin, 13 Oktober 2025, Kelas 17--22, No.~4)}}\par
    Diketahui fungsi $f(x)=\left(\dfrac{3}{x}+x^2\right)^{3/2}$. Dapatkan $f'(2)$.

    \item \textbf{[Subbab 1.4, 4.4]} {\footnotesize\textcolor{gray}{(ETS 2025/2026 -- Senin, 13 Oktober 2025, Kelas 13--16, 50--61, No.~5)}}\par
    Diketahui kurva $x^2y+4x-3y^2=2$.
    \begin{enumerate}
        \item Dapatkan $\dfrac{dy}{dx}$ dengan menggunakan turunan fungsi implisit.
        \item Dari (a) dapatkan persamaan garis singgung pada kurva tersebut di titik $(1,1)$.
    \end{enumerate}

    \item \textbf{[Subbab 4.4]} {\footnotesize\textcolor{gray}{(ETS 2025/2026 -- Selasa, 14 Oktober 2025, Kelas 5--12, 23--29, 115, 116, No.~5)}}\par
    Tentukan $\dfrac{dy}{dx}$ jika $x^2y+y^3=10$.

    \item \textbf{[Subbab 4.4]} {\footnotesize\textcolor{gray}{(ETS 2025/2026 -- Selasa, 14 Oktober 2025, Kelas 102, 108--112, No.~5)}}\par
    Diketahui $f(x)=(\sqrt{x}+2x^2+1)^5$.\\
    Dengan aturan rantai, dapatkan turunan pertama dari $f(x)$.

    \item \textbf{[Subbab 4.4]} {\footnotesize\textcolor{gray}{(ETS 2025/2026 -- Selasa, 14 Oktober 2025, Kelas 30--39, No.~5)}}\par
    Dengan diferensiasi implisit, dapatkan $\dfrac{dy}{dx}$ dari $x^3y^2-3xy^2+x=1$.

\end{enumerate}
\newpage

\end{document}
